\documentclass[12pt]{article}

\usepackage{amsmath}
\usepackage{german}
\usepackage{graphicx}
\DeclareMathOperator\sgn{sgn}

%\graphicspath{{/home/ritz/Teaching/fcn/fig/}}

%----------------- layout ----------------------------------------

\oddsidemargin=0in
\textwidth=6in
\topmargin=-0.7in
\textheight=9.5in
\baselineskip=16pt
\font\mittel=cmr17
\font\roman=cmr12
\font\klein=cmr9
\font\mgross=cmr10
\newfont{\CAP}{cmcsc10 scaled 900}
\pagestyle{empty}
\parindent=0mm

%----------------- macros ----------------------------------------
\newcommand{\conc}[1]{[{\rm C}]_{\text{#1}}}
\newcommand{\units}[1]{\text{#1}}
\newcommand{\K}{\text{K}}
\newcommand{\Na}{\text{Na}}
\newcommand{\Cl}{\text{Cl}}
\newcommand{\Leak}{\text{leak}}
\newcommand{\water}{\text{H}_2\text{O}}
\newcommand{\tderiv}[1]{\frac{\mathrm{d}#1}{\mathrm{d}t}}
%-----------------------------------------------------------------


\begin{document}

\parbox{2cm}{
\includegraphics[width=1.8cm]{../bccnlogo.pdf}
}
\parbox{11cm}{
\begin{center}
\large HUMBOLDT-UNIVERSIT"AT \hskip 0.1 cm ZU \hskip 0.1 cm BERLIN
\vskip 0.1 true cm
\mgross BERNSTEIN CENTER FOR\\ COMPUTATIONAL NEUROSCIENCE
\end{center}
}
\parbox{2cm}
{
\hfill
\includegraphics[width=1.8cm]{../hublogo.pdf}
}

\vskip 0.6 true cm

\leftline{\CAP Humboldt-Universit"at zu Berlin 
\hfill Phone: 030/2093-9110}
\leftline{\CAP Philippstr. 13 House 6
             \hfill Fax: 030/2093-6771}
\leftline{\CAP \hfill webpage: http://www.bccn-berlin.de/}


\vskip 0.8 true cm
\centerline{\bf Models of Neural Systems, WS 2010/11}
\centerline{\bf Project 1: Synfire chains}
\centerline{Project presentation and report submission: February,
15th, 2011}

\vskip 0.8 true cm
{\bf Background}
% \vskip 0.4 true cm

 A popular view in neuroscience is that the information in
 the neural systems is encoded by the number of spikes (firing rate).
 Alternatively, the information could be encoded in the precise timing
 of spikes, which would substantially increase the information capacity.
 For example, information could be processed in so-called synfire
 chains -- feedforward networks of synchronously active
 neural subpopulations. However, it has been argued that the
 background noise present in cortex would destroy all such
 temporal information. Here, you will test this hypothesis by checking if precise
 spike packets can reliably and stably propagate in model feedforward
 networks.

\vskip 0.8 true cm

{\bf Problems}

\begin{enumerate}
    \item \textit{Literature review.} Study the paper of Diesmann, Gewaltig,
        Aertsen, Nature 1999. Focus on the figures and the method
        section. What is the hypothesis? What are the main results
        and conclusions? What are the main assumptions of the model?
        What are the strengths and weaknesses of the paper?

    \item Implement the leaky integrate-and-fire model discussed in
        the computer practical (see exercise sheet 5) with the parameters given
        in Diesmann et al. (1999). Additionally, introduce the absolute
        refractory period during which the neuron does not respond to
        the incoming inputs. %If required,
        Add a relative refractory
        period as described in the paper.


    \item Model the background activity by injecting a random synaptic
        current whose values are normally distributed. Choose the mean
        and standard deviation of the current such that the neuron is
        spontaneously active with mean firing rate of 2~Hz and fires
        irregularly.

    \item The feed-forward inputs are entering the neuron model by
        means of conductance-based synapses, as discussed in computer
        practical exercise sheet 6. Choose the parameters of the alpha
        function such that the post-synaptic potential of one spike has
        the peak amplitude of 0.14 mV and time-to-peak of 1.7 ms.

    \item Model the feed-forward inputs with packets of spikes, which
        are simulated by repeatedly drawing $a_\mathrm{in}$ spike times
        from a Gaussian density with standard deviation
        $\sigma_\mathrm{in}$ and mean 0.

    \item Simulate the response of the neuron. It should respond with at
        most one spike to the pulse packet (adjust the refractory
        periods and the membrane time constant if it does not). Repeat the
        simulation multiple times and calculate the distribution of
        the output spikes. Plot and compare the input and output
        distribution.  Calculate the number of spikes $a_\mathrm{out}$ and
        standard deviation $\sigma_\mathrm{out}$ of the output spike times.
        Is the packet's propagation stable or unstable?

    \item Simulate the propagation of the pulse packet in a fully
        connected synfire chain (10 successive groups, 100 neurons
        each). Plot the raster plots (see Diesmann et al. 1999, Figure
        1).

    \item \textit{Phase-space analysis}. Calculate $a_\mathrm{out}$
        and $\sigma_\mathrm{out}$ for different pairs of
        $a_\mathrm{in}$ and  $\sigma_\mathrm{in}$. Plot the results as
        a vector field. Can you reproduce the results of Diesmann et
        al. 1999 (Figure 3)? Describe and explain possible differences.

\end{enumerate}

\vskip 0.8 true cm
{\bf Literature} 
\vskip 0.4 true cm

Markus~Diesmann, Marc-Oliver~Gewaltig and Ad~Aertsen (1999). Stable propagation
of
synchronous spiking in cortical neural networks.
\emph{Nature} 402(2):529-532.\newline

\vfill
\centerline{\CAP Contact}
\CAP

\begin{tabular}{lll}
\underline{Urs Bergmann} & Phone: 2093-8924 & Email:
urs.bergmann(at)biologie.hu-berlin.de \\
Matthias Guggenmos & & Email: matthias.guggenmos(at)bccn-berlin.de \\
Richard Kempter \hfill & Phone: 2093-8925 \hfill & Email:
r.kempter(at)biologie.hu-berlin.de \\
Bartosz Telenczuk & Phone: 2093-8838 & Email:
b.telenczuk(at)biologie.hu-berlin.de \\
\end{tabular}

\end{document}
